% Guide enseignant référent de StageCompagnon (mod_stagesynthesis et mod_stage).
% Compilation : latexmk -pdf guide-enseignant-referent.tex
\documentclass[11pt,a4paper]{article}

% Version du guide : à mettre à jour à chaque modification, avec une ligne dans l'historique.
\newcommand{\titreguide}{Guide enseignant référent}
\newcommand{\versionguide}{1.1}
\newcommand{\dateguide}{25 octobre 2026}
\newcommand{\versionstage}{0.5.1}
\newcommand{\versionsynthese}{0.2.1}

\usepackage{guide-commun}

\begin{document}

\pagedegarde{Suivre ses étudiants de toutes les promotions, valider les conventions, évaluer les stages}{70_synth_dashboard}

\tableofcontents

\section*{Historique des versions}
\addcontentsline{toc}{section}{Historique des versions}

\begin{center}
\begin{tabularx}{\linewidth}{@{}l l X@{}}
\toprule
\textbf{Version} & \textbf{Date} & \textbf{Modifications} \\
\midrule
1.0 & 16 octobre 2026 & Première version. \\
1.1 & 25 octobre 2026 & Captures d'écran refaites après le passage des listes et tableaux aux gabarits Mustache. \\
\bottomrule
\end{tabularx}
\end{center}

\newpage

% ===========================================================================
\section{Votre rôle}

En tant qu'enseignant référent, vous :
\begin{itemize}
  \item validez les demandes de convention de vos étudiants, si la scolarité
        (DEVE) l'a demandé ;
  \item évaluez leurs stages après leur auto-évaluation ;
  \item suivez leur avancement sur l'ensemble du cursus.
\end{itemize}
La validation finale et la durée retenue relèvent de la DEVE.

Vos étudiants sont répartis dans plusieurs cours de promotion
(\og Stages -- Promotion 2030\fg{}, \og Stages -- Promotion 2031\fg{}, etc.).
Le cours \og Suivi des référents de stage\fg{} les regroupe : c'est votre point
d'entrée.

% ===========================================================================
\section{Le suivi multi-promotions}

Ouvrez le cours \og Suivi des référents de stage\fg{}, puis l'activité
\og Suivi de mes étudiants\fg{}.

\subsection{Tableau de pilotage}

Une ligne par étudiant qui vous est attribué, toutes promotions confondues :
cours, année d'étude actuelle de sa promotion, thématiques obligatoires validées,
stages en attente, durée retenue et statut.

\capture{70_synth_dashboard}{Vos étudiants des promotions 2030 et 2031.}

\bouton{Voir le détail} ouvre le dossier de l'étudiant : bilan par année,
thématiques obligatoires, mobilité internationale et liste de ses stages.

\capture{72_synth_student_detail}{Dossier d'une étudiante.}

\subsection{Stages à traiter}

Le lien \bouton{Validation enseignant} ouvre la liste de travail :
\begin{itemize}
  \item en tête, les \textbf{demandes de convention à valider} ;
  \item puis les \textbf{stages à évaluer}, filtrables par étudiant, thématique
        et étape, et triables.
\end{itemize}

\capture{71_synth_entries}{Demandes de convention et stages à évaluer, toutes
promotions confondues.}

Chaque bouton (\bouton{Valider la demande}, \bouton{Évaluer}) ouvre la page
correspondante dans le cours de la promotion de l'étudiant.

\begin{astuce}
Vous ne voyez que les étudiants qui vous sont attribués. Si un étudiant manque, ou
si une promotion n'apparaît pas, contactez la DEVE : elle gère les attributions
et les promotions reliées au suivi.
\end{astuce}

% ===========================================================================
\section{Valider une demande de convention}

Vous êtes prévenu par courriel (\og Convention de stage à valider\fg{}).

\capture{26_teacher_validate}{Demande de convention à valider.}

La demande est entièrement modifiable. Vérifiez en particulier les dates, la
structure d'accueil, le tuteur et la check-list d'objectifs. Cochez
\bouton{Convention papier} si un exemplaire à signer à la main est nécessaire.

\begin{itemize}
  \item \bouton{Valider} : la demande part à la DEVE, qui édite la convention.
  \item \bouton{Refuser}, avec un motif : la demande revient à l'étudiant, qui
        reçoit le motif par courriel.
\end{itemize}

Une fois la convention éditée, vous pouvez la télécharger
(\bouton{Consulter la convention}) ; une fois signée, télécharger le PDF signé
s'il a été déposé.

% ===========================================================================
\section{Évaluer un stage}

Vous êtes prévenu par courriel quand l'étudiant s'est auto-évalué
(\og Auto-évaluation de stage à évaluer\fg{}), ou quand un import StageVet a
réuni ses évaluations (\og Stage à évaluer\fg{}).

\bouton{Évaluer} ouvre la page du stage. Elle présente le stage,
l'auto-évaluation de l'étudiant et, si elle est activée, l'évaluation du maître
de stage.

\capture{74a_teacher_eval_top}{Stage, auto-évaluation et évaluation du maître de stage.}

\begin{enumerate}
  \item \textbf{Jours de stage effectifs} : corrigez si besoin les jours
        cochés par l'étudiant, puis \bouton{Enregistrer}.
  \item \textbf{Évaluation de l'enseignant} : répondez au questionnaire, puis
        \bouton{Valider}. Le stage passe à \statut{Évalué par l'enseignant} et
        part à la DEVE pour la validation finale.
  \item Si le stage ne peut pas être retenu : \bouton{Marquer non validé}, avec
        un motif obligatoire. Le stage passe à \statut{Non validé} ; seule la
        DEVE peut le rouvrir.
\end{enumerate}

\capture{74b_teacher_eval_form}{Jours effectifs, évaluation et refus.}

% ===========================================================================
\section{Consulter le détail d'un stage}

Dans le dossier d'un étudiant, \bouton{Voir le détail} sur un stage affiche tout ce qui a été saisi : stage,
jours effectifs, convention (encadrement, stagiaire, structure d'accueil,
modalités), check-list d'objectifs et évaluations.

\capture{76_entrydetail_teacher}{Détail d'un stage.}

Depuis cette page, vous pouvez corriger la check-list d'objectifs de l'étudiant.

% ===========================================================================
\section{Depuis le cours d'une promotion}

Les mêmes actions sont disponibles dans chaque cours de promotion, page
\chemin{Validation enseignant} de l'activité \og Stages\fg{}. La liste y est
limitée à cette promotion.

\capture{73_teacher_list}{Page \og Validation enseignant\fg{} d'un cours de promotion.}

% ===========================================================================
\section{Enseignant responsable d'une thématique}

Si la DEVE vous a désigné responsable d'une thématique, la page
\chemin{Stages par thématique} liste tous les stages de cette thématique, quel
que soit le référent, avec les rapports déposés, téléchargeables un par un ou en
archive ZIP.

\capture{46_theme_stages}{Stages d'une thématique.}

% ===========================================================================
\section{Courriels que vous recevez}

\begin{center}
\begin{tabularx}{\linewidth}{@{}l X@{}}
\toprule
\textbf{Courriel} & \textbf{Action} \\
\midrule
Convention de stage à valider & Valider ou refuser la demande. \\
Auto-évaluation de stage à évaluer & Évaluer le stage. \\
Stage à évaluer (après import StageVet) & Évaluer le stage. \\
\bottomrule
\end{tabularx}
\end{center}

Chaque courriel contient un lien direct vers la page à ouvrir.

% ===========================================================================
\section{En résumé}

\begin{enumerate}
  \item Point d'entrée : cours \og Suivi des référents de stage\fg{}.
  \item Valider ou refuser les demandes de convention.
  \item Évaluer les stages auto-évalués : jours effectifs, questionnaire,
        \bouton{Valider}.
  \item Suivre l'avancement de chaque étudiant dans son dossier.
\end{enumerate}

\end{document}
